% Generated draft: requires mathematical and source review.
\documentclass[11pt]{article}
\usepackage[margin=25mm]{geometry}
\usepackage{fontspec}
\setmainfont{Latin Modern Roman}
\usepackage{amsmath,amssymb,mathtools,unicode-math}
\setmathfont{Latin Modern Math}
\usepackage{enumitem,needspace,xurl,hyperref}
\setlength{\parindent}{0pt}
\setlength{\parskip}{5pt}
\setlength{\emergencystretch}{3em}
\newcommand{\N}{\mathbb{N}}
\newcommand{\Z}{\mathbb{Z}}
\newcommand{\Q}{\mathbb{Q}}
\newcommand{\R}{\mathbb{R}}
\newcommand{\C}{\mathbb{C}}
\newcommand{\F}{\mathbb{F}}
\newcommand{\A}{\mathbb{A}}
\newcommand{\PP}{\mathbb{P}}
\newcommand{\E}{\mathbb{E}}
\newcommand{\eps}{\varepsilon}
\newcommand{\dd}{\,\mathrm{d}}
\begin{document}
\section*{UnsolvedMath ID: 800466}
\textbf{Source problem code:} catalogue entry
\subsection{Definitions and mathematical statement}
The supplied problem concerns magnetic billiards: between boundary collisions, the particle moves in a constant magnetic field along Larmor circles of radius
\[
\rho=\frac{1}{B},
\]
where \(B>0\) is the field magnitude in the normalization used by the source.

Write \(D\) for the billiard table. The intended question is whether total integrability in a strong constant magnetic field forces \(D\) to be circular. Under the usual planar interpretation of a circular table, the proposed conclusion is
\[
\exists c\in\R^2\;\exists R>0:\qquad
\partial D=\{x\in\R^2:\lVert x-c\rVert=R\}.
\]
The source does not explicitly specify the ambient space or the class of admissible tables, so this planar interpretation requires confirmation.

More formally, let \(\mathcal S\) denote the intended class of table--field pairs \((D,B)\) satisfying the source's strong-field condition, and let \(\mathrm{TI}(D,B)\) denote the intended property of total integrability. The question has the schematic form
\[
\forall (D,B)\in\mathcal S,\qquad
\mathrm{TI}(D,B)
\Longrightarrow
\exists c\in\R^2\;\exists R>0:\quad
\partial D=\{x\in\R^2:\lVert x-c\rVert=R\}.
\]
Neither \(\mathcal S\) nor \(\mathrm{TI}\) is defined by the supplied record. In particular, it gives no quantitative strong-field condition, no definition of total integrability or its phase-space scope, and no collision law, boundary regularity, or treatment of grazing trajectories. Thus this is a schematic transcription of the intended rigidity question, not a fully specified mathematical conjecture.

The background mentions convex magnetic billiards in a weak constant field, but does not explicitly impose convexity on the strong-field question. That hypothesis cannot be transferred silently. The background also distinguishes integrability from total integrability; no equivalence between these properties is supplied.
\subsection{Short English statement}
Does total integrability force a magnetic billiard table to be circular when the constant magnetic field is strong? The supplied question leaves the meanings of strong field and total integrability, as well as the admissible tables and collision conventions, unspecified. These points require human mathematical review before a precise universal statement can be fixed.
\subsection{Sources}
\begin{itemize}
\item Local UnsolvedMath registry, permanent ID \(800466\), exact source problem code \texttt{catalogue entry}; title: \emph{Billiards and Quantitative Symplectic Geometry --- Misha Bialy --- Problem 1}. The supplied registry labels the problem \emph{open}; this is a supplied status, not an independently verified mathematical assessment.
\item Albach, Barraud et al., \emph{Open Problems in Billiards and Quantitative Symplectic Geometry} (2026), source item attributed to contributor Misha Bialy, Problem 1, ``Magnetic billiards for strong magnetic fields.'' Supplied source URL: \url{https://arxiv.org/abs/2602.12896}. The record reports access on 2026-10-16 and extraction from arXiv LaTeX; these provenance details are supplied and were not independently verified here.
\item The supplied text cites the reference label \texttt{berry-robnik} for the introduction and study of magnetic billiards, \texttt{magnetic1} and \texttt{magnetic2} for algebraic approaches, \texttt{albers} for numerical experiments, and \texttt{bialy-magnetic-hopf} for the total-integrability approach for convex billiards in weak constant fields. Bibliographic entries and individual links for these labels are absent from the supplied record.
\item The background states that the circular billiard is the only known integrable example and that the total-integrability approach works well in the convex weak-field setting. These are supplied contextual literature claims, not independently verified assessments or additional hypotheses of the strong-field question.
\end{itemize}
\end{document}
